\documentclass[12 pt, letterpaper]{hmcpset}
\usepackage{hw-preamble}

\name{Jayden Thadani}
\class{MATH 177 ---  $p$-adic numbers}
\assignment{Problem set 5}
\duedate{8th October 2025}

\begin{document}

We begin by exploring quadratic extensions of $\mathbb{Q}_{p}$.

\begin{problem}[(1)]
	Suppose $f(x)$ is a polynomial of degree $2$, and let $\alpha$ be a root of $x$. We denoted the dimension of $\mathbb{Q}_{p}(\alpha)$ as a vector space over $\mathbb{Q}_{p}$ by $[\mathbb{Q}_{p}(\alpha) : \mathbb{Q}_{p}]$ and call it the \emph{degree} of the extension. For $d \in \mathbb{Q}_{p}$, denote by $\sqrt{d}$ a root of the polynomial $x^{2} - d$.
	\begin{enumerate}
		\item Show that $[\mathbb{Q}_{p}(\alpha) : \mathbb{Q}_{p}]$ equals $1$ or $2$.
		\item If $\beta$ is another root of $f(x)$, show that $\beta \in \mathbb{Q}_{p}(\alpha)$.
		\item Show that $\mathbb{Q}_{p}(\alpha) = \mathbb{Q}_{p}(\sqrt{d})$ for some $d \in \mathbb{Q}_{p}$.
		\item For $d_{1}, d_{2} \in \mathbb{Q}_{p}^{\times}$, verify that $\mathbb{Q}_{p}(\sqrt{d_{1}}) = \mathbb{Q}_{p}(\sqrt{d_{2}})$ if and only if $d_{1} d_{2}^{-1}$ is a square in $\mathbb{Q}_{p}^{\times}$.
	\end{enumerate}
\end{problem}

\pagebreak

\begin{problem}[(2)]
	Suppose $\mathbb{K}$ is an extension of $\mathbb{Q}_{p}$ of degree $2$. Show that every element of $\mathbb{K}$ is the root of a polynomial in $\mathbb{Q}_{p}[x]$ of degree $1$ or $2$.
\end{problem}

\pagebreak

\begin{problem}[(3)]
	Suppose $p > 2$.
	\begin{enumerate}
		\item Use Hensel's lemma together with a result from last week's homework to show that every non-zero $p$-adic number can be written as a product of a principal unit, a $(p - 1)$st root of unity, and a power of $p$.
		\item Let $(\mathbb{Q}_{p})^{\times} = \{ z^{2} \mid z \in \mathbb{Q}_{p} \}$. Pick $d \in \mathbb{Z}_{p}^{\times}$ for which $d \notin (\mathbb{Q}_{p}^{\times})^{2}$. Show that $\mathbb{Q}_{p}^{\times} / (\mathbb{Q}_{p}^{\times})^{2} = \{ \overline{1}, \overline{p}, \overline{d}, \overline{pd} \} \cong \mathbb{Z} / (2) \times \mathbb{Z} / (2)$.
		\item Conclude that there are exactly three extensions of $\mathbb{Q}_{p}$ of degree $2$.
		\item Something to think about but not turn in: How many extensions of $\mathbb{Q}$ of degree $2$ are there? And of $\mathbb{R}$?
	\end{enumerate}
\end{problem}

\end{document}
